% === Begin Label Data ===
% ID: 294
% 难度星级: 5.5
% 标签: 
% 备注: 
% 组卷引用次数: 0
% === End  Label Data ===

\begin{problem}{2015}{G}{江苏卷}{20}{数列}
设 $ a_1 $, $ a_2 $, $ a_3 $, $ a_4 $ 是各项为正数且公差为 $ d $ ($ d\neq 0 $) 的等差数列.

（1）证明：$ 2^{a_1} $, $ 2^{a_2} $, $ 2^{a_3} $, $ 2^{a_4} $ 依次构成等比数列；

（2）是否存在 $ a_1 $, $ d $，使得 $ a_1 $, $ a_2^2 $, $ a_3^3 $, $ a_4^4 $ 依次构成等比数列？并说明理由；

（3）是否存在 $ a_1 $, $ d $ 及正整数 $ n $, $ k $，使得 $ a_1^n $, $ a_2^{n+k} $, $ a_3^{n+2k} $, $ a_4^{n+3k} $ 依次构成等比数列？并说明理由。
\end{problem}

\begin{answer}
（1）成立；（2）不存在；（3）不存在。
\end{answer}

\begin{solutions}
（3）问．如果这样的 $ a_1 $, $ d $, $ n $, $ k $ 存在，则  
$$
a_1^n,\ a_2^{n+k},\ a_3^{n+2k},\ a_4^{n+3k} \text{ 依次构成等比数列}
$$  
$$
\Longleftrightarrow\ a_1^n,\ (a_1+d)^{n+k},\ (a_1+2d)^{n+2k},\ (a_1+3d)^{n+3k} \text{ 依次构成等比数列}
$$  
$$
\Longleftrightarrow\ n\ln a_1,\ (n+k)\ln(a_1+d),\ (n+2k)\ln(a_1+2d),\ (n+3k)\ln(a_1+3d) \text{ 依次构成等差数列}
$$  
$$
\Longleftrightarrow\ \text{存在实数 } d_1,\ b_1,\ \text{使得方程}
$$  
$$
(n+kx)\ln(a_1+dx)=d_1x+b_1
$$  
有四个根 $ 0,1,2,3 $.  

因为 $ n,k>0 $，所以当 $ x\in\{0,1,2,3\} $ 时，$ n+kx>0 $．于是方程  
$$
\ln(a_1+dx)-\frac{d_1x+b_1}{n+kx}=0
$$  
有四个不同的根 $ 0,1,2,3 $.  

设  
$$
f(x)=\ln(a_1+dx)-\frac{d_1x+b_1}{n+kx},
$$  
则  
$$
f'(x)=\frac{d}{a_1+dx}-\frac{d_1(n+kx)-k(d_1x+b_1)}{(n+kx)^2}
$$  
$$
=\frac{d(n+kx)^2-(a_1+dx)(d_1n-kb_1)}{(n+kx)^2}.
$$  

现在，注意到 $ d(n+kx)^2-(a_1+dx)(d_1n-kb_1) $ 是关于 $ x $ 的二次多项式，所以它最多只有两个不同的根，从而 $ f'(x) $ 最多只有两个不同的根，$ f(x) $ 最多只有三段单调区间，从而 $ f(x) $ 在区间 $ [0,+\infty) $ 上至多有三个相异的零点。  

因此方程 $ \ln(a_1+dx)-\dfrac{d_1x+b_1}{n+kx}=0 $ 不可能有 $ 0 $, $ 1 $, $ 2 $, $ 3 $ 四个解。  

故不存在 $ a_1 $, $ d $ 及正整数 $ n $, $ k $，使得 $ a_1^n $, $ a_2^{n+k} $, $ a_3^{n+2k} $, $ a_4^{n+3k} $ 依次构成等比数列。
\end{solutions}